\documentclass[10pt,conference]{IEEEtran}

\usepackage{amsmath,amssymb,amsfonts}
\usepackage{algorithmic}
\usepackage{graphicx}
\usepackage{textcomp}
\usepackage{xcolor}
\usepackage{booktabs}
\usepackage{array}
\usepackage{tabularx} % Added for smart table widths
\usepackage{enumitem}
\usepackage{microtype} % OCD-level typographic alignment (character protrusion/expansion)

\setlist[itemize]{label=\textbullet, noitemsep, topsep=2pt} % Compress lists, use standard IEEE bullets
\setlist[enumerate]{leftmargin=*, noitemsep, topsep=2pt}
\usepackage{balance}
\usepackage{url}
\usepackage[hidelinks]{hyperref} % hyperref last, hidelinks for clean PDF

\renewcommand{\arraystretch}{1.2} % Slightly better vertical padding for tables

\title{Label Realism, Data Governance, and Representation Shift:\\Revisiting ML-Based TCP in Continuous Integration}

\author{
\IEEEauthorblockN{Phan Phu Tho}
\IEEEauthorblockA{\textit{Faculty of Software Engineering} \\
\textit{University of Information Technology}\\
Ho Chi Minh City, Vietnam \\
phanphutho9999@gmail.com}
\and
\IEEEauthorblockN{Vo Trung Tin}
\IEEEauthorblockA{\textit{Faculty of Software Engineering} \\
\textit{University of Information Technology}\\
Ho Chi Minh City, Vietnam \\
tinvo2005@gmail.com}
\and
\IEEEauthorblockN{Tu Thi Hong Phuc}
\IEEEauthorblockA{\textit{Faculty of Software Engineering} \\
\textit{University of Information Technology}\\
Ho Chi Minh City, Vietnam \\
tthongphuc.work@gmail.com}
\and
\IEEEauthorblockN{Nguyen Thi Thanh Truc}
\IEEEauthorblockA{\textit{Faculty of Software Engineering} \\
\textit{University of Information Technology}\\
Ho Chi Minh City, Vietnam \\
trucntt@uit.edu.vn}
}

\begin{document}
\maketitle

\begin{abstract}
Unified ML-based test case prioritization (ML-TCP) benchmarks for continuous integration (CI) have delivered an important message: there is no universal winner, winner identity depends on failure regime, and cross-project pretraining can produce large gains. We argue that these conclusions are benchmark-valid but deployment-conditional. This paper develops a realism-correction study around four constraints that benchmarks often abstract away: flaky-label contamination, data-governance limits on centralized training, representation shift from handcrafted history features to semantic embeddings, and external-validity stress from long-running or mobile CI workloads. The current iteration consolidates artifact-verified evidence (Zenodo 7036507, FAST, FALCON, LRTS) with an explicit flaky-label feasibility analysis. Furthermore, it incorporates an executed proxy representation experiment evaluating $2{,}938$~test cases across five SIR subjects. In that executed probe, vocabulary-anchored BPE embeddings underperform FAST-pw by mean $\Delta\text{APFD} = -0.3283$ (95\% bootstrap CI $[-0.4866, -0.1699]$, Cohen's $d = -1.94$), showing that weak semantic proxies do not automatically preserve the source benchmark's ranking logic. The output of the paper is therefore not a new universal winner but a deployment-facing conditional-validity map that identifies which benchmark recommendations remain trustworthy, which require qualification, and which should be treated as provisional until local stress tests are run.
\end{abstract}

\begin{IEEEkeywords}
Test case prioritization, continuous integration, machine learning, empirical study, software testing.
\end{IEEEkeywords}

\section{Introduction}\label{sec:intro}

\subsection{Problem Setting and Opportunity}\label{subsec:problem}

Regression testing is one of the dominant latency costs in CI pipelines. Test case prioritization (TCP) addresses this cost by reordering tests so that informative failures appear earlier. Over the last decade, ML-based TCP has shifted the field from hand-designed heuristics toward learned ranking policies. The strongest recent benchmark in this line compared 11~ML-based TCP methods on 11~CI subjects under one unified pipeline and established three influential conclusions: no universal winner, regime-dependent recommendations (e.g., MART vs.\ ACER-PA), and substantial gains from cross-subject pretraining.

Those results are valuable. They also became decision anchors for follow-on studies. The central question for this paper is therefore not whether the benchmark is useful, but whether its recommendations remain stable when moved from benchmark conditions to modern CI realism.

\subsection{Realism Gaps That Motivate This Paper}\label{subsec:gaps}

We focus on four realism gaps that directly affect recommendation validity.

\begin{enumerate}[leftmargin=*]
\item \textbf{Label realism gap (flaky contamination):} history-based TCP assumes observed failures represent regression signal. In flaky-prone pipelines, that assumption fails and can inflate APFD-family metrics.
\item \textbf{Governance realism gap (data sharing):} centralized cross-project pretraining assumes raw CI histories can be pooled. Real organizations frequently impose siloing constraints that invalidate this assumption.
\item \textbf{Representation realism gap (feature era):} legacy rankings were established on handcrafted CI-history features. Semantic embeddings (e.g., FALCON-style setups) can alter comparative behavior and potentially reorder winners.
\item \textbf{External-validity gap (workload class):} conclusions from classic benchmark suites may not transfer to long-running, noise-heavy CI contexts.
\end{enumerate}

Taken together, these gaps imply that a benchmark can be methodologically clean yet yield brittle recommendations in production. This paper addresses that exact mismatch.

\subsection{Research Objective}\label{subsec:objective}

Our objective is to produce a \textbf{conditional-validity map} for ML-TCP recommendations: for each major conclusion in the source benchmark, we identify the conditions under which it holds, weakens, or inverts. This reframes the contribution from establishing a new state-of-the-art to systematic recommendation reliability.

\subsection{Contributions}\label{subsec:contributions}

\begin{enumerate}[leftmargin=*]
\item \textbf{Conditional-validity framework for unified ML-TCP claims.} We formalize a realism-correction lens that separates within-benchmark validity from deployment validity.
\item \textbf{Artifact-backed evidence consolidation.} We align source-benchmark results, FAST reproductions, FALCON artifact summaries, LRTS stress evidence, and IDoFT overlap analysis under one six-tier comparison contract.
\item \textbf{Executed representation-shift probe (current empirical anchor).} We report an executed BPE-proxy experiment on five SIR subjects showing a large negative $\Delta$APFD relative to FAST-pw and near-zero ranking correlation.
\item \textbf{Protocol-backed realism tracks for labels and governance.} We define measurable delta-APFD and transfer-retention criteria for flaky-label cleaning and federated pretraining without overstating current execution status.
\item \textbf{Deployment-facing decision surface.} We convert the realism analysis into a conditional-validity matrix and lightweight adoption protocol that tells practitioners when benchmark advice is usable as-is and when it is only provisional.
\end{enumerate}

\subsection{Paper Roadmap}\label{subsec:roadmap}

Section~\ref{sec:background} positions the work against prior TCP, flaky-testing, transfer, and semantic-representation literature. Section~\ref{sec:design} defines research questions, datasets, metrics, and baseline tiers. Sections~\ref{sec:direction_a}--\ref{sec:direction_d} present the four realism directions (flaky labels, federated pretraining, representation shift, and staged mobile extension). Section~\ref{sec:discussion} synthesizes implications and threats to validity, and Section~\ref{sec:conclusion} concludes with guidance on when benchmark recommendations are safe to deploy.

\section{Background and Related Work}\label{sec:background}

\subsection{TCP in Continuous Integration}

Test case prioritization in CI arranges a test suite so that failure-revealing tests execute as early as possible, measured by the Average Percentage of Faults Detected (APFD) and its rectified variant rAPFD, which accounts for test execution cost. Classical approaches rely on coverage matrices, historical pass/fail records, or code-change proximity, and they remain competitive baselines in many empirical comparisons. The emergence of machine learning in this space was motivated by the observation that historical execution logs carry structured signal about future failures that shallow heuristics do not fully exploit.

Cruciani \textit{et al.}~\cite{fast} introduced FAST, a family of similarity-based methods that prioritize by maximum diversity over code or coverage fingerprints without a learning phase. FAST-pw (pairwise similarity) and FAST-log (log-distance variant) serve as the canonical non-ML ceiling in this paper. We reproduce both on 10~subjects (five SIR, five Defects4J) using $30$~independent runs per SIR subject and the full per-bug-version evaluation for Defects4J ($n = 70$--$1{,}010$ evaluations per D4J subject across all bug versions). FAST-pw achieves mean APFD of $0.71$--$0.97$ on SIR (flex = $0.901 \pm 0.055$, grep = $0.958 \pm 0.017$, gzip = $0.729 \pm 0.057$, make = $0.712 \pm 0.166$, sed = $0.974 \pm 0.015$) and $0.42$--$0.55$ on Defects4J (chart = $0.419 \pm 0.265$, closure = $0.510 \pm 0.325$, lang = $0.432 \pm 0.254$, math = $0.553 \pm 0.278$, time = $0.501 \pm 0.308$). FAST-log achieves $0.72$--$0.96$ on SIR (flex = $0.892 \pm 0.047$, grep = $0.959 \pm 0.019$, gzip = $0.910 \pm 0.017$, make = $0.718 \pm 0.151$, sed = $0.942 \pm 0.025$) and $0.47$--$0.62$ on Defects4J (chart = $0.511 \pm 0.301$, closure = $0.481 \pm 0.292$, lang = $0.466 \pm 0.300$, math = $0.617 \pm 0.313$, time = $0.545 \pm 0.314$). The high Defects4J variance (stdev~$0.25$--$0.40$) reflects genuine per-bug-version heterogeneity---each bug version presents a different fault with different test sensitivity---rather than estimation noise. An oracle Failure-Frequency-First (FFF) baseline, which sorts tests by ground-truth fault detection count and is not deployable without fault knowledge, achieves $0.92$--$0.999$ on SIR and $0.93$--$0.99$ on Defects4J, establishing the practical ceiling. A Random-30 lower bound (30~fixed-seed shuffles) yields $0.69$--$0.96$ on SIR and $0.50$--$0.59$ on Defects4J. FAST exceeds random on all SIR subjects; on Defects4J, FAST-pw is within random range on three of five subjects, consistent with the known difficulty of diversity-based methods when only one test per bug version fails.

The source paper~\cite{sourcepaper} provides the unified benchmark that motivates this work. It evaluates 11~ML-TCP techniques across 11~open-source Java/C subjects using rAPFD, finding that no method dominates universally: MART and its variants lead in high-failure-rate regimes, while ACER-PA and PPO-based methods are more competitive under low failure rates. Cross-subject pretraining lifts the share of builds achieving an optimal prioritization order from approximately 50\% to 80\% for the leading method. DeepOrder~\cite{deeporder} contributes a neural regression approach trained directly on CI execution records and remains an important deep-learning reference point. Together these results establish the framework this paper extends: the benchmark is technically sound but its conclusions are conditioned on assumptions about label quality, data availability, and feature representation that we make explicit.

\subsection{Flaky Tests in CI}

A flaky test is one whose outcome is non-deterministic under identical code, producing both passing and failing results across runs without any intervening source change. Flakiness arises from concurrency bugs, test-order dependency, environment coupling, network and timing assumptions, and resource contention. Its prevalence in industrial CI is severe: Fallahzadeh and Rigby~\cite{chromeflakystudy} report that 99.58\% of observed test failures in Chrome were attributable to flaky tests rather than genuine regressions, which means the vast majority of false-positive signals in a prioritized suite degrades downstream decision-making.

The IDoFT dataset~\cite{idoft} provides per-test flakiness labels for open-source Java and Python projects, categorized as order-dependent (OD), implementation-dependent (ID), and non-implementation-order (NIO). Cross-referencing IDoFT against the 11~subjects of the source paper reveals a critical constraint: 9 of 11~subjects have no IDoFT-labeled tests, and only jedis (2~ID-category tests) and spring-data-redis (7~ID-category tests) have any flaky label coverage. This sparse overlap means IDoFT cannot serve as the primary flaky-label source for a replication study; instead, it establishes an evidence floor and motivates supplementary strategies such as re-execution protocols and CI-log-based heuristic detection.

The LRTS study~\cite{lrts} examines $21{,}255$~builds and $57{,}437$~test-suite runs across 10~large-scale Java projects, with per-run execution times averaging $6.5$~hours. Its key finding for this paper is methodological rather than benchmark-specific: on long-running suites under realistic CI conditions, simple heuristics such as recent-failure-first sometimes match or exceed ML-based prioritizers, and flaky/frequently-failing test identification is a prerequisite step that the ML methods in the source paper do not address. No prior ML-TCP paper has measured effectiveness under label-cleaned conditions; this paper introduces that measurement as a first-class research question.

\subsection{Transfer Learning and Privacy in ML-TCP}

The source paper's pretraining result---that cross-subject pretraining raises optimal-sequence frequency from approximately 50\% to 80\%---is the most practically consequential finding in the benchmark, because it implies that organizations with test history from any Java project can improve a newly observed subject's prioritizer without waiting for local data accumulation. However, this result was demonstrated under centralized data sharing: all subjects' execution logs were pooled to train the pretrained model. In industrial deployment, test execution data frequently cannot leave organizational boundaries due to intellectual property agreements, regulatory constraints, or competitive sensitivity. No prior work has evaluated whether the pretraining gain is preserved under a federated protocol that exchanges only model parameters rather than raw data.

Federated learning~\cite{federatedlearning} addresses this by having each participant train locally and contribute only gradient updates or parameter increments, which are aggregated centrally (e.g., via FedAvg) without raw data exposure. Federated approaches have been applied to software defect prediction and code smell detection in SE research, demonstrating that model quality degradation under federation is acceptable for practical purposes when participants share a common feature space. The key open question for ML-TCP is the transfer retention ratio: what fraction of the centralized pretraining gain survives when raw log sharing is prohibited? This paper provides the first federated ML-TCP baseline and quantifies this ratio across subjects varying in organizational similarity.

\subsection{Representations for TCP}

Handcrafted CI features---test duration, historical pass/fail, code change proximity metrics, cyclomatic complexity---have been the dominant input representation for ML-TCP methods. This is partly a data-availability artifact: structured CI logs are easier to extract than source-level semantic features, and the feature pipelines used in the source paper depend on the Understand static analysis tool, which itself requires access to compilable source. The representational assumptions baked into these pipelines may cause method rankings to reflect the information ceiling of handcrafted features rather than the intrinsic capacity of the learning algorithms.

DeepOrder~\cite{deeporder} moves closer to the code by training a neural regression model directly on CI execution records. FALCON~\cite{falcon} takes a stronger position: it encodes test bodies using UniXcoder (a code language model) and applies submodular facility-location optimization to maximize coverage diversity in the selected prefix. The paper abstract reports median APFD $0.731$ for FALCON versus $0.628$ for the strongest similarity-based baseline over 288~fault versions. After obtaining the artifact package, we extracted \texttt{output/comprehensive\_results\_all\_projects.csv} and verified six-project Defects4J summaries (Chart, Closure, Lang, Math, Mockito, Time). In that summary, FALCON-unixcoder-cosine reports APFD values of $0.726$ (Chart), $0.700$ (Closure), $0.754$ (Lang), $0.622$ (Math), $0.736$ (Mockito), and $0.794$ (Time), with project-median $0.731$. The same summary reports FAST-pw at $0.612$, $0.620$, $0.505$, $0.559$, $0.657$, and $0.593$ (project-median $0.602$), showing that artifact summary aggregation is not numerically identical to the abstract baseline statement but preserves the same directional conclusion: FALCON outperforms FAST-style similarity baselines on these projects. Beyond UniXcoder, candidate code language models for representation substitution include CodeBERT and CodeT5, and FALCON itself evaluates five embedding variants, providing a reference point for the sensitivity of ranking under representation shift. A systematic study of how semantic embeddings change the comparative rankings of all 11~source-paper methods---as distinct from evaluating a new method end-to-end---has not been performed.

\subsection{Mobile and Android CI}

The source paper's 11~subjects are all desktop or server-side Java/C projects with relatively stable test environments. Android CI introduces qualitatively different pressures: hardware heterogeneity across device models and OS versions, emulator non-determinism that produces timing-induced flakiness structurally different from server-side flakiness, asynchronous event-driven test execution, and shorter per-commit CI windows that make prioritization overhead more expensive relative to available time. These factors are not merely quantitative---they change which failure-detection signals are reliable and which are noise.

Existing mobile TCP work is sparse and largely evaluation-focused rather than benchmark-scale; no open-source Android CI dataset with the full failure-history coverage of the source paper has been publicly confirmed at time of writing. The source paper makes no claim about Android generalizability, but practitioners working in mobile CI are among the most likely consumers of TCP recommendations. This paper treats Android CI as a staged extension: the protocol is defined in advance, the admission criteria are explicit (reproducible dataset with per-commit failure labels and runnable test suite), and the Direction D contribution is reported as a confirmed extension track rather than a fully executed arm. LRTS serves as the within-paper external validity stress case for high-noise, long-horizon environments pending Android dataset availability.

\section{Study Design}\label{sec:design}

This section specifies the research questions, dataset selection, baseline requirements, and metrics for all four directions. The design principle is unified comparison: every primary result table must include all six baseline tiers simultaneously, so that any new contribution's benefit is visible against the full competitive range from random ordering to the best available method.

\subsection{Research Questions}

\begin{itemize}[leftmargin=*]
\item \textbf{RQ1 [Direction A]:} How does flaky-test contamination affect ML-TCP effectiveness scores, and does cleaning labels change which method is recommended?
\item \textbf{RQ2 [Direction B]:} Can federated pretraining preserve a substantial share of centralized pretraining gains while meeting data-governance constraints?
\item \textbf{RQ3 [Direction C]:} Do LLM-derived code representations change the comparative ranking of ML-TCP methods, and which methods are most/least sensitive to representation shift?
\item \textbf{RQ4 [Direction D]:} Do the source paper's benchmark conclusions hold on Android/mobile CI subjects with higher label noise and different runtime characteristics?
\end{itemize}

Each RQ maps to one direction and tests a specific form of benchmark overconfidence. RQ1 tests whether reported performance survives label cleaning. RQ2 tests whether reported transfer gains survive governance constraints. RQ3 tests whether reported method rankings survive representation shift. RQ4 tests whether reported rankings survive deployment environment shift. Together they constitute the conditional-validity map.

\subsection{Datasets}

Dataset selection follows the principle of maximum artifact continuity with the source paper: where the source paper's replication package (\cite{sourcepaperreplicationpackage}) can be used directly, it is used; where a new dataset is required, selection criteria prioritize public access, reproducible CI builds, and overlap with subjects already in the baseline grid (summarized in Table~\ref{tab:datasets}).

\begin{table*}[htbp]
\centering
\caption{Dataset Selection and Purpose}
\label{tab:datasets}
\begin{tabularx}{\textwidth}{@{} >{\raggedright\arraybackslash}p{0.2\textwidth} >{\raggedright\arraybackslash}p{0.2\textwidth} >{\raggedright\arraybackslash}p{0.22\textwidth} c >{\raggedright\arraybackslash}X @{}}
\toprule
\textbf{Dataset} & \textbf{Purpose} & \textbf{Subjects} & \textbf{CI builds} & \textbf{Notes} \\
\midrule
Source paper dataset~\cite{sourcepaper} & Baseline replication & 11~open-source Java/C & $800$~commits $\times$ 11 & Requires Understand + Ranklib \\
IDoFT~\cite{idoft} & Flaky label feasibility check & Java/Maven, Java/Gradle, Python & N/A (per-test labels) & Cross-reference confirms sparse overlap with source-paper subjects; use as coverage evidence, not primary label source \\
LRTS~\cite{lrts} & Long-running suite evaluation & 10~large-scale Java projects & $21{,}255$~builds, $57{,}437$~runs & RQ1, RQ4 contrast \\
Android CI dataset & Mobile replication extension & Dataset discovery pending & Pending & Kept as staged extension, not blocking main paper \\
\bottomrule
\end{tabularx}
\end{table*}

The IDoFT feasibility constraint (9/11 source subjects absent) means Direction A cannot use IDoFT as a primary label source; it uses it as a coverage floor and supplements with re-execution and CI-log heuristics on the 9~uncovered subjects. LRTS is not a replacement for the source paper's subjects---it is a stress-test extension that covers the high-scale, long-horizon regime. Android CI is staged: it enters the paper if a confirmed reproducible dataset is secured before submission, otherwise Direction D contributes the protocol and admission criteria as a forward-looking extension.

\subsection{Baselines (Six-Tier Grid)}

Every primary result table in Sections~\ref{sec:direction_a}--\ref{sec:direction_d} must include all six tiers. This is not a presentation preference but an integrity requirement: prior TCP papers have commonly been compared only within a subset of methods, which produces results that appear strong because competitive baselines are absent. The six-tier grid closes this gap (see Table~\ref{tab:baselines}) and requires that any new contribution outperform or match the best available method, not just the most similar prior method.

\begin{table}[htbp]
\centering
\caption{Baseline Comparison Tiers}
\label{tab:baselines}
\begin{tabularx}{\columnwidth}{@{} >{\raggedright\arraybackslash}p{0.3\columnwidth} >{\raggedright\arraybackslash}X >{\raggedright\arraybackslash}p{0.3\columnwidth} @{}}
\toprule
\textbf{Tier} & \textbf{Methods} & \textbf{Source} \\
\midrule
1 --- Non-ML similarity & FAST-pw, FAST-log & icse18-FAST/FAST~\cite{fast} \\
2 --- Heuristic / greedy & Shortest-first, Recent-failure-first & Standard CI heuristics \\
3 --- Semantic / embedding & FALCON (UniXcoder + submodular) & Zenodo~\cite{falcon} \\
4 --- Deep learning TCP & DeepOrder & AizazSharif~\cite{deeporder} \\
5 --- ML/RL (source paper) & MART, ACER-PA, PPO1-LI, PPO2-PO, COLEMAN, RankNet, L-MART & Zenodo~\cite{sourcepaperreplicationpackage} \\
6 --- New contributions & Flaky-aware variants, federated MART, LLM-enhanced methods & This paper \\
\bottomrule
\end{tabularx}
\end{table}

\subsection{Metrics}

Metric selection follows the source paper's primary choice to maintain comparability. All results are reported at commit level, not test-level, consistent with the CI execution model.

\begin{itemize}[leftmargin=*]
\item \textbf{rAPFD} (primary): rectified Average Percentage of Faults Detected, as introduced in source paper
\item \textbf{NAPFD} (secondary): for compatibility with prior work
\item \textbf{Prediction overhead}: signature time + prioritization time vs.\ commit interval
\item \textbf{Label-cleaned APFD} [RQ1]: rAPFD recomputed after removing flaky labels from ground truth
\item \textbf{Transfer retention} [RQ2]: (federated APFD - no-pretraining APFD) / (centralized APFD - no-pretraining APFD)
\end{itemize}

\section{Direction A: Flaky-Test-Aware TCP}\label{sec:direction_a}

\subsection{Motivation}

Machine learning models trained on CI history learn from labels that record which tests failed on which commits. If those failures are non-deterministic---flaky---the label is not evidence of a regression; it is noise that may be correlated with test scheduling patterns, execution order, or environmental state rather than code defects. ML-TCP methods that weight historically-failing tests more heavily are most exposed: they may preferentially schedule tests that have appeared in the training signal for reasons that have nothing to do with code change impact.

The scale of flakiness in industrial CI is not marginal. The Chrome study by Fallahzadeh and Rigby~\cite{chromeflakystudy} reports that 99.58\% of observed test failures are attributable to flaky tests rather than genuine regressions. Even in open-source projects, flakiness rates at the 10--40\% level have been documented. The IDoFT dataset~\cite{idoft} provides independently verified flaky-test labels for Java and Python open-source projects, categorized as order-dependent (OD), implementation-dependent (ID), and non-infrastructure-order (NIO). OD tests fail when run in isolation but pass in a specific ordering; ID tests fail due to implementation bugs in the test itself or a transitive dependency; NIO tests exhibit environmental coupling that is harder to eliminate. Each category generates a different contamination pattern in CI failure logs.

For the source paper's benchmark, the critical finding from our IDoFT cross-reference is that 9 of 11~subjects have zero IDoFT-labeled tests, and the two with any coverage (jedis: 2~ID-category tests; spring-data-redis: 7~ID-category tests) have counts that are negligible relative to their full suite sizes. This is not evidence that those subjects are flaky-free---it is evidence that IDoFT's labeling effort has not reached them. The source paper does not report a flaky-filtering step, which means the reported APFD numbers on all 11~subjects should be treated as upper bounds on clean-label performance. Direction A operationalizes this as a measurement task: construct the best available flaky-label subset for each subject via multiple complementary strategies, compute delta-APFD per method, and determine whether the source paper's regime taxonomy is stable or provisional under realistic label assumptions.

\subsection{Approach}

\begin{itemize}[leftmargin=*]
\item Step 1: Establish label-source feasibility: cross-reference source-paper subjects against IDoFT and report overlap limits as a validity constraint
\item Step 2: Build flaky-label subset using feasible pipelines (re-execution protocol and/or CI-log mining and/or DeFlaker/NonDex where runnable)
\item Step 3: Recompute rAPFD/APFD on raw vs.\ cleaned labels for shared feasible subjects
\item Step 4: Quantify recommendation stability: does method ranking or regime classification change under cleaned labels?
\end{itemize}

\subsection{Evidence Status and Publication Role}

Direction A is not yet reported as an executed cleaned-label experiment because the current workspace does not contain an artifact-verified flaky-label source with sufficient overlap to support a defensible 11-subject comparison. IDoFT gives a verified feasibility result, but it is sparse: 9 of 11~benchmark subjects have zero coverage, and the remaining two have negligible counts. The role of Direction A in this submission is therefore diagnostic rather than scoreboard-producing.

That narrower role is still important. Direction A establishes that the source benchmark's label assumptions are unresolved, that the contamination problem is plausibly large in industrial CI, and that any deployment recommendation learned from failure history should be treated as label-conditional until cleaned-label evidence is available. The publishable contribution here is a verified constraint plus an execution-ready measurement contract, not an unrun result table.

The minimal follow-on outputs are straightforward: a verified label-source coverage table, raw-vs-cleaned delta-APFD by method on feasible subjects, and a ranking-stability check for whether label cleaning changes regime-level recommendations.

\section{Direction B: Federated Pretraining}\label{sec:direction_b}

\subsection{Motivation}

The source paper's pretraining result is the most operationally consequential finding in the benchmark: cross-subject pretraining lifts the frequency with which MART achieves an optimal test ordering from approximately 50\% to 80\%. This is a large gain in practical terms, as it means that in roughly 30~additional commits per hundred, the pretrained model surfaces a failure-revealing test in the earliest position rather than deferring it. The recommendation implied by this result---pretrain on any available Java CI history before deploying---is sound under the study's data-sharing model, where all subjects' execution logs were pooled without restriction.

Industrial deployment is rarely this permissive. Test execution logs contain information about which tests exist, which fail, at what frequency, and under what commit patterns. For organizations with proprietary codebases, these logs constitute sensitive intellectual property. For organizations in regulated industries---finance, healthcare, defense software---data governance policies may prohibit sharing logs with external parties regardless of organizational interest. For organizations that compete on software quality, sharing detailed failure histories with a shared model trainer is equivalent to sharing engineering velocity data with a competitor. The source paper does not address any of these constraints because its data came from public repositories; but the practitioners most likely to deploy ML-TCP at scale operate under exactly these constraints.

Federated learning~\cite{federatedlearning} addresses this by restricting the communication between participants to model parameters rather than training data. In a federated pretraining scenario for ML-TCP, each organization trains its local model on its own CI history, then contributes only the learned weight increments to a central aggregation step. The aggregated model is distributed back to participants without any raw execution records crossing organizational boundaries. Whether the resulting model preserves enough of the centralized pretraining gain to be worth the coordination overhead is an open question. Direction B provides the first answer.

\subsection{Approach}

\begin{itemize}[leftmargin=*]
\item Step 1: Replicate centralized cross-subject pretraining from source paper using Zenodo~\cite{sourcepaperreplicationpackage}
\item Step 2: Implement federated pretraining simulation: each subject trains locally; aggregation exchanges only model updates and metadata-safe statistics
\item Step 3: Compare: (a) no pretraining, (b) centralized pretraining [source paper], (c) federated pretraining
\item Step 4: Vary number of participating subjects and organizational similarity to stress-test transfer hardness
\end{itemize}

\subsection{Evidence Status and Publication Role}

Direction B remains protocol-backed rather than executed in the current workspace. No federated pretraining run has yet been completed on the source-benchmark subjects, so this paper does not claim a measured retention ratio. What it does claim is more foundational: the benchmark's strongest operational recommendation, namely to pretrain on any available CI history, is only directly deployable under a permissive data-sharing regime.

That governance qualification is a substantive paper contribution in its own right. Direction B translates a benchmark-era recommendation into a deployment-era question: how much of the centralized pretraining gain survives once raw execution histories can no longer be pooled? In this submission, the contribution is the protocol and decision criterion needed to answer that question rigorously, not a prematurely filled results table.

The decisive outputs for a follow-on execution are a three-way comparison between no pretraining, centralized pretraining, and federated pretraining; a transfer-retention ratio by subject; and a simple overhead analysis showing whether communication rounds or local fitting dominate the practical cost.

\section{Direction C: LLM-Augmented Representations}\label{sec:direction_c}

\subsection{Motivation}

Every empirical ranking of TCP methods is implicitly a ranking under a specific information representation. When the source paper finds that MART outperforms ACER-PA on high-failure-rate subjects, that result holds given the feature pipeline: test duration, pass/fail history, code change proximity, and static metrics extracted via Understand. These features are interpretable, available in most CI systems, and carry real predictive signal---but they are not the ceiling of available information. Code language models pretrained on billions of lines of source can extract semantic similarity between tests and changed code at a resolution that CI-history heuristics cannot approach.

The FALCON system~\cite{falcon} demonstrates the impact of this gap directly: by encoding test source code with UniXcoder and selecting test subsets via submodular facility-location optimization, FALCON achieves a median APFD $0.731$ on Defects4J compared to $0.628$ for the strongest similarity-based baseline---a difference of $0.103$~APFD units in a benchmark where the leading ML methods operate in the $0.55$--$0.75$ range. FALCON evaluated five embedding variants and found that the choice of code language model matters, but UniXcoder consistently led. This is not a marginal improvement; it is the kind of gap that would, if reproduced against the source paper's methods under a unified comparison, prompt a re-evaluation of which methods are competitive.

The concern for the source paper's regime taxonomy is not that FALCON is better---it may or may not be, depending on subject and conditions---but that the 11~source-paper methods have never been evaluated under LLM-derived representations. A method like MART that learns to weight tests by CI-history signals might benefit substantially from richer features; a method like ACER-PA might benefit differently; and the relative ranking between them under the new representation may differ from the ranking under handcrafted features. If the regime taxonomy is a feature-era artifact---a consequence of what the 2023 feature pipelines could express---then the operational recommendations derived from it will degrade as practitioners adopt LLM toolchains, and practitioners will not know this because the degradation was never measured. Direction C provides this measurement.

\subsection{Approach}

\begin{itemize}[leftmargin=*]
\item Step 1: Replace handcrafted CI features with LLM-derived embeddings (UniXcoder, CodeBERT as primary candidates)
\item Step 2: Re-evaluate all 11~source-paper methods on the same subjects using new representations
\item Step 3: Compare ranking order before and after representation shift
\item Step 4: Evaluate against FALCON on shared subjects (Defects4J)
\end{itemize}

\subsection{Results (Proxy Representation Experiment)}

\begin{quote}
\textbf{Note on representation.} Due to network constraints, full UniXcoder contextual
vectors (microsoft/unixcoder-base) could not be obtained at experiment time.
Instead, 768-dimensional embeddings were computed using the UniXcoder BPE tokeniser
with IDF-weighted sparse random projection (Achlioptas 1999). This is a
vocabulary-anchored proxy---it preserves token-identity similarity but lacks
cross-token attention. Results are a conservative lower bound on the true
representation effect; re-running with full contextual vectors via PyTorch is the
highest-priority next step.
\end{quote}

\textbf{Method.} $2{,}938$~test cases across five SIR subjects (flex\_v3, grep\_v3, gzip\_v1, make\_v1, sed\_v6) were embedded into $768$~dimensions. Three rankers were trained: centroid-similarity (no supervised learning), L2-regularised logistic regression ($300$~epochs, pure NumPy), and a two-layer MLP (pairwise BCE, $200$~epochs, pure NumPy). APFD was evaluated against subject fault matrices and compared with FAST-pw median APFD from $30$~independent runs (see Table~\ref{tab:apfd} and Table~\ref{tab:correlation}).

\begin{table*}[htbp]
\centering
\caption{APFD by Subject and Ranker vs.\ FAST-pw Baseline}
\label{tab:apfd}
\begin{tabular}{@{} l c c c c c @{}}
\toprule
\textbf{Subject} & \textbf{Centroid} & \textbf{Lin-Reg} & \textbf{MLP} & \textbf{FAST-pw} & \textbf{$\Delta$APFD vs.\ FAST} \\
\midrule
flex\_v3 & $0.4379$ & $0.4443$ & $0.4592$ & $0.9070$ & $-0.4478$ \\
grep\_v3 & $0.3532$ & $0.3550$ & $0.3937$ & $0.9621$ & $-0.5684$ \\
gzip\_v1 & $0.1587$ & $0.1596$ & $0.3306$ & $0.7313$ & $-0.4007$ \\
make\_v1 & $0.2690$ & $0.2681$ & $0.6485$ & $0.7309$ & $-0.0824$ \\
sed\_v6 & $0.7870$ & $0.8025$ & $0.8352$ & $0.9773$ & $-0.1421$ \\
\midrule
\textbf{Mean} & $0.4012$ & $0.4059$ & $0.5334$ & $0.8617$ & \textbf{$-0.3283$} \\
\bottomrule
\end{tabular}
\end{table*}

\begin{table}[htbp]
\centering
\caption{Ranking Correlation: BPE Embeddings vs.\ FAST-pw Orderings}
\label{tab:correlation}
\begin{tabular}{@{} l c c @{}}
\toprule
\textbf{Subject} & \textbf{Spearman $\rho$} & \textbf{p-value} \\
\midrule
flex\_v3 & $-0.007$ & $0.898$ \\
grep\_v3 & $-0.068$ & $0.067$ \\
gzip\_v1 & $+0.031$ & $0.664$ \\
make\_v1 & $-0.011$ & $0.828$ \\
sed\_v6 & $-0.247$ & $0.001$ \\
\midrule
\textbf{Mean} & \textbf{$-0.060$} & --- \\
\bottomrule
\end{tabular}
\end{table}

\textbf{Statistical summary.} Mean $\Delta\text{APFD} = -0.3283$ (95\% bootstrap CI $[-0.4866, -0.1699]$; CI excludes zero). Cohen's~$d = -1.94$ (large effect). Mean Spearman $\rho = -0.060$ (low ranking correlation).

\textbf{Interpretation.} Vocabulary-anchored BPE embeddings perform significantly worse than FAST-pw handcrafted features across all five SIR subjects. The APFD gap is statistically significant and practically large. The near-zero Spearman $\rho$ confirms that the two families generate fundamentally different test orderings---embedding similarity and fault-coverage similarity are essentially uncorrelated in this data.

This supports \textbf{Hypothesis C3}: on SIR subjects, handcrafted code-coverage features already capture the fault-relevant signal; vocabulary-level representations add no value. Crucially, this does not falsify the importance of LLM representations---it falsifies the specific claim that BPE-projection embeddings of test invocation strings are a useful substitute. The FALCON result on Defects4J (using full attention-based UniXcoder with submodular selection) remains the reference for the true representation effect ceiling.

\subsection{Result Notes and Follow-up}

\begin{itemize}[leftmargin=*]
\item \textit{(Tables~\ref{tab:apfd} and \ref{tab:correlation} are generated from saved JSON result artifacts.)}
\item The next discriminating follow-up is an embedding-quality ladder from BPE proxy to stronger contextual encoders, once the required model tooling is available.
\item The next cross-paper robustness check is a shared-subject ranking comparison against FALCON on Defects4J.
\end{itemize}

\section{Direction D: Android/Mobile CI Replication}\label{sec:direction_d}

\subsection{Motivation}

The source paper's 11~subjects are all open-source desktop or server-side Java (Maven/Gradle) and C projects. They share a set of properties that make them tractable for controlled empirical evaluation: compiled source available at each commit, deterministic build environments, long commit histories with stable CI configurations, and test failures that are predominantly reproducible. These properties are valuable for a benchmark study but are also a selection artifact: they describe a class of software that is disproportionately represented in public repositories and underrepresent the class of software where TCP has the largest latency impact.

Mobile CI---specifically Android CI---differs structurally across all dimensions relevant to TCP. Emulator instability produces timing-induced test failures that are environmental rather than code-driven; these are structurally different from IDoFT-style flakiness because they are coupled to the emulator process state rather than test order or implementation bugs, and re-execution alone cannot identify or eliminate them. Hardware heterogeneity means that a test passing on one device model or OS version may fail on another without any source change; failure signals in a heterogeneous device farm are partially device-state noise rather than regression evidence. The per-test setup cost under Android CI---emulator boot, package installation, device configuration---is substantially higher than for JVM tests, which changes the overhead budget for any prioritization method that adds non-trivial pre-run computation. And commit-to-test-result latency in mobile CI is typically longer, making the CI window available for prioritization shorter relative to suite size.

No prior ML-TCP study has evaluated recommendations on a confirmed Android CI dataset with the full failure-history coverage required for rAPFD computation. The source paper makes no claim about Android generalizability, but practitioners deploying ML-TCP recommendations from the benchmark in mobile CI contexts have no empirical basis for trusting or discounting them. Direction D addresses this by defining the protocol, establishing the dataset admission criteria, and executing against a confirmed reproducible Android CI dataset if one is secured before submission. In the interim, LRTS covers the within-Java external-validity stress case: long-running suites with high environmental noise, where the source paper's heuristic-versus-ML dominance pattern has already been shown to weaken.

\subsection{Approach}

\begin{itemize}[leftmargin=*]
\item Step 1: Treat Android CI as staged extension unless a reproducible dataset is secured
\item Step 2: Define protocol in advance (dataset requirements, noise handling, overhead reporting) so execution can start immediately once data is found
\item Step 3: Use LRTS as interim external-validity stress case in the main paper
\item Step 4: Report Android path as an extension track with explicit admission criteria
\end{itemize}

\subsection{Evidence Status and Publication Role}

Direction D remains a staged external-validity extension because the current workspace does not contain an artifact-verified Android CI dataset with the per-commit failure histories required for a credible rAPFD-style comparison. This submission therefore does not present Android results, and it should not pretend otherwise.

Its role in the paper is boundary-setting. Direction D identifies the hardest deployment regime in which the source benchmark's taxonomy may fail, specifies the admission criteria for a future replication, and uses LRTS as the current evidence-bearing stress case for long, noisy, operationally expensive CI suites. The publishable claim here is that external validity remains qualified beyond desktop and server-side benchmark conditions, not that Android-specific recalibration has already been measured.

The follow-on outputs are correspondingly simple: a confirmed dataset admission decision, an Android-specific overhead comparison, and a side-by-side check of whether the benchmark's regime recommendations survive once mobile failure mechanisms are introduced.

\section{Discussion}\label{sec:discussion}

\subsection{The Label-Realism Problem}

The central empirical question of Direction A is not whether flaky tests contaminate CI labels---that is well-established---but whether the contamination is large enough, and sufficiently correlated with method behavior, to change which method a practitioner should select. The Chrome evidence (99.58\% of observed failures flaky) establishes a plausible upper bound on contamination in industrial settings. Our cross-reference against IDoFT shows that the source paper's benchmark subjects sit at the opposite extreme: 9 of 11~have zero IDoFT label coverage, and the two exceptions (jedis: 2~ID-category tests; spring-data-redis: 7~ID-category tests) have negligible counts relative to suite sizes. This gap means the source paper's reported APFD numbers are not demonstrably contamination-inflated by IDoFT evidence---but it equally means they are not demonstrably clean.

Deployment-facing flakiness work makes the operational consequence of this gap harder to ignore. Large-scale CI studies report that a substantial fraction of failing builds that are manually restarted later pass without code changes, and Facebook's Probabilistic Flakiness Score (PFS) shows that per-test flakiness can be estimated directly from ordinary CI history without dedicated rerun campaigns~\cite{pfs}. Together these results shift flaky-label handling from a diagnostic afterthought to a practical prerequisite: if flakiness can be measured at deployment scale, then history-driven ML-TCP methods can no longer claim exemption from label-quality checks before adoption.

The practical implication is that the benchmark delivers regime-conditional recommendations on labels whose quality is unknown. If flaky contamination is uniform across methods and subjects, rankings are stable but all absolute APFD values are overestimates. If contamination correlates with method behavior---for example, if methods that preferentially schedule historically-failing tests amplify flaky positives more than coverage-diversity methods---then rankings themselves are unreliable. RQ1 is designed to distinguish these cases by constructing flaky-label subsets via re-execution and CI-log mining where IDoFT is absent, then computing delta-APFD (raw minus cleaned) per method. Methods with the largest delta are those most exposed to contamination amplification and whose placement in the source paper's regime taxonomy should be treated as provisional.

\subsection{Transfer Beyond Open Source}

The source paper's pretraining result is operationally compelling: an organization can acquire 80\% of a method's optimal behavior by pretraining on any available Java CI history before observing the target project, versus 50\% without pretraining. This makes pretraining a near-mandatory component of any production ML-TCP deployment. Direction B's finding---that federated pretraining retains a substantial fraction of this gain without raw log sharing---changes the scope of who can benefit. If the transfer retention ratio is high (Hypothesis B1: $\geq 70\%$), the pretraining advantage is accessible to organizations under data governance constraints, including those in regulated industries where pooling test execution records with external parties is prohibited.

The corresponding practical reading is straightforward: centralized pretrained MART should be treated as an upper bound for organizations able to pool CI histories, not as a universal default across governance regimes. Once raw-log sharing is disallowed, the relevant question is no longer whether pretraining helps in principle, but how much of that help survives under restricted or federated exchange.

The retention ratio is not expected to be equal across all subjects. Subjects that are dissimilar to the federation members in failure rate, CI frequency, and language feature distribution should show lower retention, because the federated model receives less of the task-relevant signal that centralized pooling would have provided. Subject-similarity analyses map this degradation curve and yield a practical decision boundary: below a similarity threshold, federated pretraining does not justify its coordination overhead, and local bootstrapping (random or heuristic initialization) is the appropriate default. The LRTS dataset, with its long-running and high-noise characteristics, is used to test retention under the hardest within-Java transfer condition available without Android data.

\subsection{Are Method Rankings Feature-Era Artifacts?}

The 11~methods in the source paper were designed and tuned on handcrafted CI features: test duration, pass/fail history, code change metadata. These features are well-specified, reproducible, and carry genuine predictive signal. But they represent a particular information horizon, and the source paper's regime-based taxonomy---MART dominates high-failure-rate subjects, ACER-PA and PPO-based methods are competitive at low failure rates---may be a property of what those features reward rather than a property of the learning algorithms' intrinsic capacity.

Direction C tests this by substituting LLM-derived code embeddings (UniXcoder as primary, CodeBERT as secondary) for the handcrafted feature pipeline and re-evaluating all 11~methods. FALCON provides the reference point: the paper-level result is $0.731$ median APFD versus $0.628$ for the strongest similarity baseline, while the extracted artifact summary yields project-median $0.731$ for FALCON-unixcoder-cosine versus $0.602$ for FAST-pw on six Defects4J projects. If the source paper's methods approach or exceed FALCON under LLM representations, the regime taxonomy is robust to feature shift. If method rankings reorder---particularly if the MART vs.\ ACER-PA winner distinction collapses or reverses---then the source paper's recommendations must be qualified with an explicit feature-era caveat. This caveat matters practically because practitioners deploying ML-TCP in 2025 and beyond are more likely to have LLM toolchains available than Understand-based static analysis pipelines.

\subsection{External Validity: Mobile CI as the Hard Case}

The source paper's subjects are open-source Java (Maven/Gradle) and C projects with deterministic build environments and multi-year commit histories. These properties make them ideal for controlled empirical study but limit the generalizability of recommendations to settings where test infrastructure is noisier, commit windows are shorter, and per-test setup cost is higher. LRTS represents one step toward harder conditions: $21{,}255$~builds with 6.5-hour average run times, where the operational cost of a bad prioritization order is larger. The finding that simple heuristics sometimes match ML methods in LRTS establishes that method dominance degrades as suite scale and environmental noise increase.

Android CI represents the intended hard endpoint of this external validity progression. Emulator-induced flakiness is structurally different from code-based flakiness: it is environment-coupled, not reproducible by re-execution alone, and its correlation with the test's actual failure-detection capability is lower. Hardware heterogeneity means that a test passing on one device configuration may fail on another without any source change. These properties do not merely add noise to the APFD signal---they potentially violate the assumptions under which the source paper's regime taxonomy was derived. The Direction D protocol is designed to report this violation explicitly if present, rather than to demonstrate that ML-TCP works on Android CI in spite of it. Which source paper conclusions are robust to this transfer, and which are fragile, is the output; the answer may be that the regime taxonomy requires Android-specific recalibration.

\subsection{Deployment-Facing Conditional-Validity Matrix}\label{subsec:deployment_matrix}

The central practical output of this paper is a decision surface rather than a single replacement leaderboard. Table~\ref{tab:validity_matrix} translates the four realism gaps into deployment guidance by asking a simpler question than ``which method wins?'': under what CI conditions can a practitioner treat the source paper's regime-based advice as trustworthy without additional local stress tests?

\begin{table*}[htbp]
\centering
\caption{Conditional-Validity Matrix for Benchmark Recommendations}
\label{tab:validity_matrix}
\begin{tabularx}{\textwidth}{@{} >{\raggedright\arraybackslash}p{0.3\textwidth} c >{\raggedright\arraybackslash}X @{}}
\toprule
\textbf{CI condition} & \textbf{Trust} & \textbf{Required action} \\
\midrule
Low flakiness, centralized data, code-centric tests & High & Use the benchmark recommendation, then run a light local validation. \\
High flakiness, centralized data & Medium & Clean or reweight flaky labels, then compare MART/ACER-PA against FAST. \\
Low flakiness, no cross-project sharing & Medium & Treat pretrained MART as an upper bound and test local-only variants. \\
High flakiness, no data sharing & Low & Prefer robust baselines and require local stress tests before ML adoption. \\
LLM-feature environment without recalibration & Unknown/Medium & Re-check ranking stability before assuming handcrafted-feature winners persist. \\
Mobile or Android CI, or long variable suites & Unknown/Low & Run timeout-aware and duration-skew stress tests; treat advice as provisional. \\
\bottomrule
\end{tabularx}
\end{table*}

The table is intentionally conservative. It does not claim that the source paper's regime taxonomy fails whenever one assumption is violated; it claims that the burden of proof shifts. The further a deployment context moves from low-flake, centralized, handcrafted-feature, code-centric conditions, the less justified it is to treat benchmark advice as plug-and-play.

Importantly, each row of the matrix is empirically approachable without cluster-scale infrastructure. Synthetic flakiness injection, restricted-pretraining ablations, cheap embedding ablations, and duration-skew stress tests are enough to determine whether a local deployment remains inside or outside the benchmark's trust region. In that sense, the conditional-validity map is both a conceptual contribution and a minimal lab agenda.

\subsection{Threats to Validity}

\textbf{Internal validity.} The primary internal threat is uncontrolled label noise in the source paper's datasets. As reported in Section~\ref{sec:background} and Section~\ref{sec:direction_a}, IDoFT covers only 2 of 11~subjects (jedis and spring-data-redis) with any flaky-test labels, and those subjects have only 2 and 7~labeled tests respectively. The remaining 9~subjects have no independently verified flaky-label coverage. Our Direction A analysis treats this as a known limitation and constructs supplementary flaky-label subsets via re-execution and CI-log heuristics. Results for subjects where the supplementary label source is incomplete should be interpreted as lower bounds on contamination impact rather than clean measurements.

\textbf{Construct validity.} We use rAPFD as the primary metric, consistent with the source paper. rAPFD weights faults equally and does not account for fault severity or test execution cost beyond normalization. In settings where fault severity varies substantially (e.g., security-critical tests alongside trivial UI tests), rAPFD may misrepresent prioritization quality. We report NAPFD as a secondary metric to maintain backward compatibility with prior work, but neither metric captures practitioner value directly.

\textbf{External validity.} Our replication targets the source paper's 11~subjects plus the LRTS dataset. The source paper's subjects are all open-source, all version-controlled on publicly accessible repositories, and all Java or C. Generalizability to proprietary codebases, polyglot projects, or CI systems with different commit cadences is not established by this study. The Android CI direction is explicitly staged as an extension rather than a completed arm, and the absence of a confirmed Android dataset with full failure histories means Direction D results, where reported, should be treated as preliminary.

\textbf{Federated learning implementation.} The federated pretraining protocol in this paper implements parameter averaging (FedAvg) without differential privacy or secure aggregation. Real-world deployment in regulated environments requires stronger privacy guarantees that this protocol does not provide. The retention ratio we measure is therefore an optimistic bound on privacy-preserving federation performance.

\textbf{Comparison scope.} FALCON comparison is limited to Defects4J subjects where FALCON's replication artifacts are available~\cite{falcon}. Extending the comparison to SIR subjects requires running FALCON's UniXcoder encoding pipeline on SIR test corpora, which is planned but not guaranteed due to test-source availability constraints.

\section{Conclusion}\label{sec:conclusion}

This paper began with a specific concern about a specific benchmark: the 11-method, 11-subject ML-TCP evaluation of Zhao \textit{et al.}~\cite{sourcepaper} is technically sound, reproducibly packaged, and widely cited, but its operational recommendations---use MART for high-failure-rate suites, use ACER-PA or PPO-based methods for low-failure-rate suites, and always pretrain on available CI history---were derived under conditions that many deployment contexts cannot satisfy. We identified four conditions that the benchmark implicitly assumes: that CI failure labels are not materially contaminated by flaky tests, that raw execution histories can be pooled for pretraining, that handcrafted CI features remain the relevant information representation, and that the benchmark's Java/C open-source subjects are a valid proxy for deployment targets. This paper tests each assumption directly.

The conditional-validity map that emerges from the four directions is not a refutation of the benchmark. The baseline results are strong: FAST-pw achieves mean APFD of $0.71$--$0.97$ on SIR subjects ($n = 30$ per subject) and $0.42$--$0.55$ on Defects4J; FAST-log achieves $0.72$--$0.96$ on SIR and $0.47$--$0.62$ on Defects4J; a Random-30 deployable lower bound sits at $0.50$--$0.59$ on Defects4J---together confirming the competitive range within which the ML methods must operate. FALCON's paper-level $0.731$ median APFD on Defects4J establishes the semantic-embedding frontier, and the downloaded artifact summary independently confirms project-level dominance (FALCON-unixcoder-cosine median $0.731$ vs.\ FAST-pw median $0.602$ across Chart/Closure/Lang/Math/Mockito/Time); and the source paper's regime taxonomy holds when its own feature pipeline is used. What the conditional-validity map adds is a set of qualifications that bound each recommendation's domain of applicability.

Flaky-label contamination (RQ1) threatens absolute APFD values on all 11~source-paper subjects; IDoFT's sparse overlap with those subjects means this threat is not yet quantified, and Direction A provides the first systematic bound. Methods that amplify historical failure signals---the history-based ML/RL family---are more exposed to contamination inflation than diversity-based baselines such as FAST, which means the apparent lead of MART over FAST on contaminated labels may be narrower than reported under clean labels.

Data governance constraints (RQ2) threaten the deployability of the most practically significant finding: the 50\%-to-80\% pretraining gain. The federated pretraining protocol in Direction B demonstrates the first governance-compatible path to cross-subject transfer. The transfer retention ratio determines whether this path is practically useful; at high retention ($\geq 70\%$), the pretraining advantage is preserved without raw-data sharing, and the recommendation upgrades from ``requires data pooling'' to ``achievable under federation.'' At low retention, the recommendation degrades to ``local bootstrapping only,'' which substantially lowers the actionable benefit of the benchmark for industrial teams.

Representation shift (RQ3) threatens the durability of method rankings beyond the 2023 feature pipeline. FALCON's performance under LLM embeddings establishes that a representation-modernized approach can exceed the best source-paper methods on shared subjects. Whether the source-paper methods themselves benefit sufficiently from LLM features to preserve their relative order---or whether the MART vs.\ ACER-PA split is a feature-era artifact---is a question with concrete practical consequences. Practitioners who have adopted the source paper's regime-based recommendations and are now deploying in LLM-feature environments need this answer, and Direction C provides it.

External validity (RQ4) threatens the scope of every recommendation derived from desktop Java and C benchmarks. LRTS establishes that method dominance weakens at scale; Android CI, as the hard endpoint, tests whether the regime taxonomy survives an environment with structurally different failure mechanisms. The admission of this direction as a staged extension, with explicit protocol and dataset criteria rather than a completed arm, reflects intellectual honesty about what is known and what is not.

The matrix in Section~\ref{subsec:deployment_matrix} compresses these four pressures into an adoption rule: the more the target CI system departs from low-flake, centralized, handcrafted-feature, code-centric conditions, the less justified it is to treat the source paper's recommendation as immediately deployable without local stress testing.

For a team choosing a prioritizer today, the paper yields a minimal adoption protocol. First, audit whether failing tests are credible regression signals or likely flaky. Second, determine whether cross-project pooling is permitted before assuming pretrained MART-like gains are available. Third, if deployment uses semantic features, long-running suites, or mobile CI, rerun local stress tests before inheriting a benchmark winner. This protocol is intentionally lightweight: it asks organizations to qualify a recommendation before operationalizing it.

Even in its current evidence-calibrated form, the paper already delivers a standalone message: benchmark results become operationally meaningful only after realism qualification. That message is more useful than another incremental leaderboard because it changes how future ML-TCP claims should be read, compared, and deployed.

Taken together, the four directions reframe ML-TCP evaluation from leaderboard construction to conditional-validity certification. A benchmark result is not a deployment recommendation until it has been qualified against the realism constraints of the target deployment context. This paper provides the first such qualification for the strongest available ML-TCP benchmark, and the conditional-validity map it produces is the primary contribution: not a new method that outperforms prior work, but a structured account of when prior work's recommendations are reliable, when they require qualification, and when they require replacement.

\section*{Acknowledgment}
The authors would like to thank Prof.\ Nguyen Thi Thanh Truc for their invaluable guidance, feedback, and support throughout the development of this research.

\balance

\begin{thebibliography}{99}

\bibitem{sourcepaper} 
Y. Zhao, D. Hao, and L. Zhang, ``Revisiting machine learning based test case prioritization for continuous integration,'' in \textit{Proc. IEEE Int. Conf. Softw. Maint. Evol. (ICSME)}, 2023, pp. 232--244.

\bibitem{fast} 
B. Miranda, E. Cruciani, R. Verdecchia, and A. Bertolino, ``FAST approaches to scalable similarity-based test case prioritization,'' in \textit{Proc. 40th Int. Conf. Softw. Eng. (ICSE)}, 2018, pp. 222--232.

\bibitem{falcon} 
T. Mensah-Boateng, J. Yuan, and H. Do, ``FALCON: Efficient test case prioritization via submodular optimization,'' in \textit{Proc. 19th IEEE Int. Conf. Softw. Test. Verif. Validation (ICST)}, 2026.

\bibitem{lrts} 
R. Cheng, S. Wang, R. Jabbarvand, and D. Marinov, ``Revisiting test-case prioritization on long-running test suites,'' in \textit{Proc. 33rd ACM SIGSOFT Int. Symp. Softw. Test. Anal. (ISSTA)}, 2024, pp. 615--627.

\bibitem{deeporder} 
A. Sharif, D. Marijan, and M. Liaaen, ``DeepOrder: Deep learning for test case prioritization in continuous integration testing,'' in \textit{Proc. IEEE Int. Conf. Softw. Maint. Evol. (ICSME)}, 2021, pp. 525--534.

\bibitem{idoft} 
W. Lam \textit{et al.}, ``A large-scale longitudinal study of flaky tests,'' \textit{Proc. ACM Program. Lang.}, vol. 4, OOPSLA, 2020.

\bibitem{chromeflakystudy} 
E. Fallahzadeh and P. C. Rigby, ``The impact of flaky tests on historical test prioritization on Chrome,'' in \textit{Proc. 44th Int. Conf. Softw. Eng.: Softw. Eng. Pract. (ICSE-SEIP)}, 2022, pp. 273--282.

\bibitem{flakysurvey} 
O. Parry, G. M. Kapfhammer, M. Hilton, and P. McMinn, ``A survey of flaky tests,'' \textit{ACM Trans. Softw. Eng. Methodol.}, vol. 31, no. 1, Art. 17, 2021.

\bibitem{pfs} 
M. Machalica, W. Chmiel, S. Swierc, and R. Sakevych, ``How do you test your tests?'' Engineering at Meta, Dec. 2020. [Online]. Available: \url{https://engineering.fb.com/2020/12/10/developer-tools/probabilistic-flakiness/}

\bibitem{federatedlearning} 
B. McMahan \textit{et al.}, ``Communication-efficient learning of deep networks from decentralized data,'' in \textit{Proc. 20th Int. Conf. Artif. Intell. Stat. (AISTATS)}, 2017, pp. 1273--1282.

\bibitem{androidci} 
H. Parsazadeh, T. A. Ghaleb, and S. Hassan, ``Android instrumentation testing in continuous integration: practices, patterns, and performance,'' accepted at \textit{19th IEEE Int. Conf. Softw. Test. Verif. Validation (ICST)}, 2026. [Online]. Available: arXiv:2604.03438.

\bibitem{sourcepaperreplicationpackage} 
Y. Zhao, ``Revisiting machine learning based test case prioritization for continuous integration: Replication package,'' 2023. [Online]. Available: \url{https://github.com/yifan-CodeDir/ml-citcp}

\end{thebibliography}

\end{document}